\cotefile{157-1}
\scanpage{18}
\mspage{15}
% archive sheet number: 17
\margin{\deleted{\unsure{18}}}
Par exemple, considérons le dérivateur
\[
\mathrm{HOTAB}
\]
dont la cat. \unsure{fondamentale} est $\mathrm{Hotab}$, la \inserted{sous-cat. pleine de la} cat. dérivée
\deleted{des complexes de \uline{chaînes}} \inserted{$\mathrm D(\mathrm{Ab})$} \deleted{de} $\mathrm{Ab}$, \deleted{\unsure{soit}} formée des complexes \deleted{\unsure{qui sont}} à coh. nulle en degrés $>0$.
On peut décrire $\mathrm{HOTAB}$ par \deleted{\unsure{la|les}} catégorie de modèles $\underline{\mathrm{Hom}}(\unsure{\underline{\Delta}^{\circ}|\Delta^{\circ}},\mathrm{Ab})$ des groupes abéliens $1/2$ simpliciaux, avec comme équivalences \deleted{faibles} les quasi-isom. au sens habituel. Le ``\uline{foncteur d'abélianisation}'' \deleted{\unsure{de passage au}} \inserted{classique}
\[
\mathrm{Hot}\longrightarrow\mathrm{Hotab}
\]
se prolonge en un morphisme de dérivateurs
\[
\mathrm{HOT}\longrightarrow\mathrm{HOTAB}
\]
La question de sa compatibilité avec $f_*$ ne semble pas tautologique. Elle \deleted{\unsure{revient}} impliquerait p.ex. que si
\diagram{\begin{tikzcd}[ampersand replacement=\&]
X \arrow[d,"f"'] \& X' \arrow[l,"g"'] \arrow[d,"f'"]\\
Y \& Y' \arrow[l,"{\unsure{j|\delta}}"]
\end{tikzcd}}
est un diagramme cartésien dans $\mathcal M=\underline{\mathrm{Hom}}(\Delta^\circ,\mathrm{Ens})$,
avec $f$ une fibration (\unsure{en restant} dans le cadre des esp. top., \unsure{une|en} \inserted{p.ex.} fibration loc. triviale) alors le \unsure{carré} dans $\mathrm{Hotab}$
% slanted
\margin{$\unsure{\longleftrightarrow|\longleftarrow}$ C'est faux déjà si \deleted{\unsure{le}} $Y$ est l'espace ponctuel, \unsure{puisque|très} \unsure{visiblement} le $\mathbb R\Gamma_!$ ne transforme pas produit en \unsure{une|en|ou} \unsure{somme} ! Ce foncteur \unsure{n'est guère|n'est pas un} ``exact à \unsure{gauche|g.}'' (comme morphisme de dérivateurs)}
\diagram{\begin{tikzcd}[ampersand replacement=\&]
\mathbb R\Gamma_!(X)\arrow[d] \& \mathbb R\Gamma_!(X')\arrow[l]\arrow[d]\\
\mathbb R\Gamma_!(Y) \& \mathbb R\Gamma_!(Y')\arrow[l]
\end{tikzcd}}
est \inserted{hot-}cartésien. \deleted{\illegible{2}} Et de $\unsure{=|\Leftarrow}$, \unsure{pour|par} une hot-\deleted{\unsure{so}}\unsure{somme amalgamée} (carré \inserted{hot-}cocartésien dans $\mathrm{Hot}$).
\scanpage{19}
\mspage{16}
% archive sheet number: 18
Si c'est vrai, et pas connu de plus, ce serait quand $\unsure{=|\simeq}$ scandaleux !

En procédant comme pour les morphismes de chgt de base pour les $Rf_*$, on trouve en tous cas des morphismes canoniques, pour $f:I\longrightarrow J$
\[
u_J f_*^{\mathbb D}\longrightarrow f_*^{\mathbb D'}u_I
\]
et de $=$
\[
f_!^{\mathbb D'}u_I\longrightarrow u_J f_!^{\mathbb D}\,,
\]
et la question est de savoir si c'est des iso.
Questions analogues pour les $f^!$, pour imm. \unsure{fermées|fermée}, (\deleted{\unsure{et}} pour $\unsure{f^?|f_!}$, pour $f$ imm. \unsure{ouverte|ouvertes}). Comme la notion de morphisme de dérivateurs devrait quand $=$ ``coller'' à la notion bien connue de morphisme \inserted{(triangulé)} de cat. triangulées, donc être ``exact'' dans un sens convenable, elle montre bien $*$ que l'on doit exiger la commutation des \deleted{\unsure{fonct}} morphismes de dérivateurs aux $f_*$ et aux $f_!$.
% slanted
\margin{Non ! Le cas des cat. triangulées est ``misleading'', car, pour celles-ci, ``exactitude à g. et à dr.'' \unsure{reviennent} au même.}
% slanted
\margin{\uline{NB} On conçoit qu'un foncteur entre dérivateurs puisse être exact à g. ou exact à dr., ou les deux !}

\unsure{J|Je} devrais demander à Thomason s'il a regardé, pour un morphisme
\[
\mathcal M\longrightarrow\mathcal M'
\]
entre cat. de modèles, qui \deleted{soit} \deleted{\unsure{exact}} disons, \uline{exact à gauche}, et ``compatible avec les \inserted{équivalences,} fibrations (et cofibrations'' (???)).
\margin{Notation $\mathbb{DER}$ pour la 2-catégorie des \uline{dérivateurs}}

\figure{\begin{tikzpicture}\draw[line width=.6pt] (0,.20) -- (.16,0);\end{tikzpicture}}
\scanpage{20}
\mspage{17}
% archive sheet number: 19
\uline{Dualité} Si $\mathbb D$ est un dérivateur, on va définir un dérivateur \uuline{dual} $\mathbb D^\circ$ par
\[
\mathbb D^\circ(I)\overset{\text{déf}}{=}\mathbb D(I^\circ)^\circ
\]
Les coefficients sur $I$ \deleted{\unsure{pour}} de type $\mathbb D^\circ$, sont les \unsure{coefficient|coefficients} sur la cat. opp. $I^\circ$ à $I$ de type $\mathbb D$, et en tant que cat. c'est la cat. \uline{opposée} de celle des coeff. de type $\mathbb D$ sur $I^\circ$. Pour définir la loi 2-fonctorielle, on note que pour $I,J$ dans $\mathrm{Diag}$ on a
\diagram{\begin{tikzcd}[ampersand replacement=\&,column sep=large]
\underline{\mathrm{Hom}}(I,J)^\circ \arrow[r,"{\mathbb D_{I,J}}"] \arrow[d,"\text{dualité}","{\rotatebox{90}{$\scriptstyle\sim$}}"'] \& \underline{\mathrm{Hom}}(\mathbb D(J),\mathbb D(I))\arrow[d,"\text{dualité}","{\rotatebox{90}{$\scriptstyle\sim$}}"']\\
\underline{\mathrm{Hom}}(I^\circ,J^\circ)\arrow[r,"\Theta"] \& \underline{\mathrm{Hom}}(\mathbb D(J)^\circ,\mathbb D(I)^\circ)^\circ
\end{tikzcd}}
où $\Theta$ est l'unique foncteur rendant le diagramme commutatif. Alors $\Theta^\circ$ est un foncteur
\[
\underline{\mathrm{Hom}}(I^\circ,J^\circ)^\circ\xrightarrow{\unsure{D|\mathbb D}}\underline{\mathrm{Hom}}(\mathbb D(J)^\circ,\mathbb D(I)^\circ),
\]
c'est ça par définition $\mathbb D^\circ_{I,J}$.

Si $\mathbb D$ est défini par $(\mathcal M,\Sigma)$, alors \inserted{$\mathbb D^\circ$ par} $(\mathcal M^\circ,\Sigma^\circ)$, puisque
\[
\deleted{\unsure{\mathrm{Hom}}}\,(\mathcal M^\circ)_I\overset{\text{déf}}{=}
\underline{\mathrm{Hom}}(I^\circ,\mathcal M^\circ)\simeq
\underline{\mathrm{Hom}}(\underset{\substack{\shortparallel\\(I^\circ)^\circ}}{I},\mathcal M)^\circ
\simeq\deleted{\unsure{\mathbb D(I)}}(\mathcal M_{I^\circ})^\circ
\]
et $\Sigma^\circ_I$ est déduit de $\Sigma_{I^\circ}$, \deleted{\unsure{par}} d'où
\[
\big((\mathcal M^\circ)_I\,\Sigma^\circ_I\big)^{-1}\simeq
\big(\mathcal M_{I^\circ}\,\Sigma_{I^\circ}^{-1}\big)^\circ=\mathbb D(I^\circ)^\circ
\]
\scanpage{21}
\mspage{18}
% archive sheet number: 20
\uline{Philosophie} Chaque fois qu'on a \deleted{deux \unsure{dérivateurs}} \inserted{un foncteur}
\[
\mathcal A\xrightarrow{u_0}\mathcal A'
\]
entre deux catégories qui peuvent s'interpréter \inserted{naturellement} comme les catégories fondamentales de deux dérivateurs $\mathbb D,\mathbb D'$, il se pose la question d'interpréter $u_0$ comme provenant d'un morphisme de dérivateurs
\[
\mathbb D\xrightarrow{u}\mathbb D'
\]
Celui-ci est, en pratique, tjrs déterminé canoniquement, à iso. canonique près. (Mais je doute pourtant que le foncteur ``restriction aux cat. fondamentales''
\[
\underline{\mathrm{Hom}}(\mathbb D,\mathbb D')\longrightarrow\underline{\mathrm{Hom}}(\mathcal A,\mathcal A')
\]
soit pleinement fidèle \ldots)

On va appliquer cette philosophie aux foncteurs $f^*$, $f_!$, $f_*$ (et éventuellement $f^!$, $f^?$), pour $f:I\longrightarrow J$, donc pour des foncteurs entre $\mathbb D(I)$ et $\mathbb D(J)$. \unsure{Cela|Ceci} nous amène donc à regarder $\mathbb D(I)$ et $\mathbb D(J)$ comme des cat. fondamentales de dérivateurs convenables.

\uline{Dérivateur }\uuline{induit} par un dérivateur $\mathbb D$, sur une catégorie d'indices $I\in\unsure{\mathrm{Dia}|\mathrm{Di}}$. C'est un dérivateur $\mathbb D_I$ défini par
\[
\mathbb D_I(J)=\mathbb D(I\times J)
\]
Le foncteur associé à $\mathbb D_I$, pour deux cat $J,J'\in\mathrm{Dia}$
\scanpage{22}
\mspage{19}
% archive sheet number: 21
\diagram{\begin{tikzcd}[ampersand replacement=\&,column sep=small,row sep=small]
{\underline{\mathrm{Hom}}(J,J')^\circ}\arrow[rr,"{(\mathbb D_I)_{J,J'}}"] \& \& {\underline{\mathrm{Hom}}(\mathbb D_I(J'),\mathbb D_I(J))}\arrow[d,equal]\\
\& \& {\underline{\mathrm{Hom}}(\mathbb D(I\times J'),\mathbb D(I\times J))}\arrow[ddr,out=-70,in=90]\\
\text{est le composé} \& \& \\
{\underline{\mathrm{Hom}}(J,J')^\circ}\arrow[r] \& {\underline{\mathrm{Hom}}(I\times J,I\times J')^\circ}\arrow[r,"{\mathbb D_{I\times J,I\times J'}}"] \& {\deleted{\unsure{\underline{\mathrm{Hom}}}}\unsure{)}} \& {}
\end{tikzcd}}
où la première flèche est \deleted{\unsure{déduite de}} \inserted{le} foncteur
\[
\deleted{\unsure{g}}\, f\longmapsto\mathrm{id}_I\times f
\]
En d'autres termes, $\mathbb D_I$ est le composé
\[
\mathrm{Diag}^\circ
\xrightarrow[\displaystyle J\mapsto\underbrace{I\times J}_{=J_I}]{\phantom{\mathbb D}}
\mathrm{Diag}^\circ\xrightarrow[\mathbb D]{}\mathrm{Cat}
\]
On a
\[
f^*_{\mathbb D_I}=(\mathrm{Id}\times f)^*_{\mathbb D}
\]
Si les $f_*$ (resp. les $f_!$) existent \unsure{pour} $\mathbb D$, ils existent pour $\mathbb D_I$. (De $\unsure{=|\simeq}$ pour les espèces de multiadjoints.) Comme $J\longrightarrow I\times J$ transforme immersion fermée (resp. ouverte) en \unsure{itou|idem}, l'axiome $\mathrm{Der}_2$ (\deleted{\unsure{exp}}plus précisément, chacun de \unsure{ces|ses} ingrédients, cf \unsure{annotations|annot.} marginales p. 7) est stable par passage de $\mathbb D$ à $\mathbb D_I$. De même pour $\mathrm{Der}_3$. (Pour chaque \unsure{nouvel|nouvelle} axiome qu'on dégagera, il faudra vérifier la stabilité.)
\scanpage{23}
\mspage{20}
% archive sheet number: 22
Soit
\[
f:I\longrightarrow I'
\]
Considérons entre $\mathbb D(I),\mathbb D(I')$ les foncteurs
\diagram{\begin{tikzcd}[ampersand replacement=\&,column sep=huge]
\mathbb D(I)\arrow[r,bend left=32,"f_!"]\arrow[r,bend right=32,"f_*"'] \& \mathbb D(I')\arrow[l,"f^*"']
\end{tikzcd}}
(sous réserve d'existence de $f_*$ resp. $f_!$). Peut on les interpréter en termes de morphismes de dérivateurs
\diagram{\begin{tikzcd}[ampersand replacement=\&,column sep=large]
\mathbb D_I\arrow[r,bend left=32,"\bar f_!"]\arrow[r,bend right=32,"\bar f_*"'] \& \mathbb D_{I'}\arrow[l,"\bar f^*"'] \& {?}
\end{tikzcd}}
Donc il faut définir, pour \deleted{\unsure{tous|tout}} $J$ \inserted{dans $\mathrm{Diag}$}, des foncteurs correspondants entre $\mathbb D_I(\deleted{\unsure{I}}J)$ et $\mathbb D_{I'}(J)$, i.e. entre $\mathbb D(I\times J)$ et $\mathbb D(I'\times J)$. Or on a bien
\[
\deleted{\unsure{\mathrm{id}\times f}} \quad f\times\mathrm{id}_J:I\times J\longrightarrow I\times J'
\]
et on posera (sous réserve d'existence \unsure{des} deux derniers)
\[
\left\{\begin{aligned}
\bar f^{*}_{\deleted{\unsure{*}}}(J)&=(\unsure{F|f}\times\mathrm{id}_J)^*_{\mathbb D}\\
\bar f_!(J)&=(f\times\mathrm{id}_J)^{\mathbb D}_!\\
\bar f_*(J)&=(f\unsure{\times|\vee}\mathrm{id}_J)^{\mathbb D}_*
\end{aligned}\right.
\]
\scanpage{24}
\mspage{21}
% archive sheet number: 23
Autrement dit, on peut considérer
\[
\underset{\substack{\rotatebox{-90}{$\scriptstyle\in$}\\\mathrm{Diag}}}{I}
\longmapsto
\underset{\substack{\rotatebox{-90}{$\scriptstyle\in$}\\\mathrm{DER}}}{\mathbb D_I}
\]
comme fonctoriel \deleted{\unsure{en}} en $I$ de trois façons \inserted{(\unsure{une contr.}, les deux autres ``co'')}, suivant qu'on utilise les $f^*$, les $f_!$ (s'ils existent) ou les $f_*$ (s'ils existent). On désignera par
\[
f^{\mathbb D}_!,\quad f^*_{\mathbb D},\quad f^{\mathbb D}_*
\]
indifféremment les foncteurs entre $\mathbb D(I)$ et $\mathbb D(I')$, ou (de façon plus fine) entre $\mathbb D_I$ et $\mathbb D_{I'}$, dont les précédents se déduisent en passant aux catégories fondamentales correspondantes. De $\unsure{=|\simeq}$ pour $f^!$ ou pour $f^?$ sous réserve d'existence. P. ex. si on suppose que les $f^!$ existent pour $\mathbb D$, quel que soit $f$ imm. fermée, alors on trouve \inserted{pour \deleted{\unsure{tout}} tel $f$} un morphisme de dérivateurs
\[
\mathbb D_I\xleftarrow{f^!_{\mathbb D}}\mathbb D_{I'}.
\]

Il s'impose de regarder si tous \unsure{les|ces} morphismes de dérivateurs sont exacts, i.e. si pour
\[
g:J\longrightarrow J'
\]
dans $\mathrm{Diag}$, ils commutent à $g_!$, $\unsure{g_*|J_*}$.
\scanpage{25}
\mspage{23}
% archive sheet number: 24
Pour expliciter cette question, on est amené à regarder le diagramme de catégories
\diagram{\begin{tikzcd}[ampersand replacement=\&,row sep=large,column sep=huge]
I\times J\arrow[r,"{G=\mathrm{id}_I\times g}"]\arrow[d,"{F=f\times\mathrm{id}_J}"'] \& I\times J'\arrow[d,"{F'=f\times\mathrm{id}_{J'}}"]\\
I'\times J\arrow[r,"{G'=\mathrm{id}_{I'}\times g}"] \& I'\times J'{,}
\end{tikzcd}}
qui est d'ailleurs un carré cartésien dans $(\mathrm{Cat})$.
La commutation des $f^*_{\mathbb D}$ aux $g_*$ s'explicite alors par la commutativité \unsure{du|des}
\diagram{\begin{tikzcd}[ampersand replacement=\&,column sep=large,row sep=large]
\mathbb D(I\times J)\arrow[r,"G_*",""{name=top,below}] \& \mathbb D(I\times J')\\
\mathbb D(I'\times J)\arrow[u,"F^*"]\arrow[r,"G'_*"] \& \mathbb D(I'\times J')\arrow[u,"F'^*"',""{name=right,left}]\arrow[from=right,to=top,dashed,bend left=20,"\text{iso?}"]
\end{tikzcd}}
Cette commutativité sera assurée par des axiomes ultérieurs sur les dérivateurs. (Cf plus bas \inserted{page \unsure{59|39|29}}) \unsure{Alors|Ainsi}, $f^*_{\mathbb D}$ est exact à droite, et d'autre part, il doit être \inserted{aussi} exact à gauche. Donc $f^*_{\mathbb D}$ \inserted{$(:\mathbb D_{I'}\unsure{-|\to}\mathbb D_I)$} \uline{est exact}.

\deleted{\unsure{On a}} Considérons $f_*^{\mathbb D}$, cette fois c'est la commutativité du
\diagram{\begin{tikzcd}[ampersand replacement=\&,column sep=large]
\mathbb D(I\times J)\arrow[d,"F_*"'] \& {\overset{G_*}{\unsure{\longrightarrow|\mathord{-}}}} \& \mathbb D(I\times J')\arrow[d,"F'_*"]\\
\mathbb D(I'\times J)\arrow[rr,"G'_*"] \& \& \mathbb D(I'\times J'){,}
\end{tikzcd}}
c'est trivial \deleted{\unsure{par}} (à iso can. près). Ainsi, on
\scanpage{26}
\mspage{24}
% archive sheet number: 25
trouve que $f^{\mathbb D}_*:\mathbb D_I\longrightarrow\mathbb D_{I'}$ est \uline{exact} \uline{à gauche}, et \unsure{dualmt} $f^{\mathbb D}_!$ est \uline{exact à droite}. Mais est-ce que $f^{\mathbb D}_*$ est aussi exact à \uline{droite} ? Cela signifie donc comm. du
\diagram{\begin{tikzcd}[ampersand replacement=\&,column sep=large,row sep=large]
\mathbb D(I\times J)\arrow[r,"G_!"]\arrow[d,"F_*"] \& \mathbb D(I\times J')\arrow[d,"F'_*",""{name=side,left}]\\
\mathbb D\unsure{(|\,}I'\times J)\arrow[r,"G'_!",""{name=base,above}] \& \mathbb D(I'\times J')\arrow[from=base,to=side,bend left=20]\mathpunct{.}
\end{tikzcd}}
On arrive bien à définir
\[
G'_!F_*\xrightarrow{?}F'_*G_!
\]
ou ce qui revient au $\unsure{=|\simeq}$ par adjonction
\[
F_*\xrightarrow{?}G'^*F'_*G_!
\]
en utilisant le morphisme de chgt de base
\[
\unsure{\sim|\,}G'^*F'_*\xrightarrow{\substack{\alpha\\\sim}}F_*G^*
\]
qui est \uline{iso} (on l'a vu), d'où
\diagram{\begin{tikzcd}[ampersand replacement=\&,row sep=small]
\& G'^{*}F'_*G_!\arrow[dd,"\alpha*G_!","{\rotatebox{90}{$\scriptstyle\sim$}}"']\\
F_*\arrow[ur,"?"]\arrow[dr,"F_* *\beta"'] \& \\
\& F_*\unsure{G_1^*|G^*}G_{\unsure{!|*}}
\end{tikzcd}}
où $\beta:\mathrm{id}\longrightarrow G^*G_!$ \deleted{$\unsure{G|f}$}\unsure{)} est le morphisme d'adjonction, d'où le morphisme voulu. (On peut aussi chercher $F'^*(G'_!F_*]\longrightarrow G_!$, et utiliser l'\deleted{\unsure{exacti}} commutation $G_!F^*\xrightarrow{\sim}F'^*G'_!$, i.e. ex. à \uline{droite} \unsure{du|des} \deleted{$\unsure{\mathbb D_I\longrightarrow}$} $f^*_{\mathbb D}$)
\scanpage{27}
\mspage{25}
% archive sheet number: 26
Mais ce morphisme est-il un \uline{iso}morphisme ? Je pense que dans le cas additif ça peut être raisonnable, mais dans le cas non additif ça n'a aucune raison d'être vrai. On voit en tous cas que si pour tout $f:I\longrightarrow I'$ dans $\mathrm{Diag}$, $f_*^{\mathbb D}:\mathbb D_I\longrightarrow\mathbb D_{I'}$ est (non seulement exact à g., mais aussi) exact à droite, alors pour tout $g:J\longrightarrow J'$ dans $\mathrm{Diag}$, $g_!^{\mathbb D}$ est (non seulement exact à droite, mais même) exact à gauche.

Prenons le dérivateur des préfaisceaux d'ens. sur des $I$ variables, sa cat. fondamentale est $\mathrm{Ens}$. \inserted{c'est un dérivateur non pointé} Considérons $I\longrightarrow I'$ une imm. fermée \deleted{\unsure{et}} \inserted{non \deleted{\illegible{1}} surjective} $f_{\unsure{*|!}}:\mathbb D(I)\longrightarrow\mathbb D(I')$. Si $f_*^{\mathbb D}$ était \unsure{exacte|exact} à droite, \deleted{il} \inserted{$f_*$} devrait transformer objet initial de $\mathbb D(I)$ en objet initial de $\mathbb D(\unsure{I'|I})$, sauf erreur. (Cf plus bas) Or il n'en est rien, puisque pour \unsure{tout} $F\in\mathrm{Im}(f_*^{\mathbb D}:\mathbb D(I)\longrightarrow\mathbb D(I'))$ la restriction de $F$ à $I'\setminus I=U$ est le faisceau final $e_U=(i\longmapsto\{e\})$, et non le faisceau initial $\varnothing_U=(i\longmapsto\varnothing)$.

\unsure{Donc on prouvera ceci|On va prouver ceci} :

\uline{Prop} : \deleted{\unsure{Si un}} $\mathbb D$ est un dérivateur à g. (à dr.), alors les produits finis (sommes finies) existent dans les $\mathbb D(I)$. Tout \deleted{\unsure{f}} morphisme \deleted{de} $\varphi:\mathbb D\unsure{\longrightarrow|-}\mathbb D'$ de dérivateurs qui est exact à g (à dr.) commute auxdits produits (sommes), \deleted{\unsure{ceci}} pour chaque $I$ donnant $\varphi(I):\mathbb D(I)\longrightarrow\mathbb D'(I)$. En particulier il transforme objet final (initial) en \unsure{itou|idem}
\scanpage{28}
\mspage{26}
% archive sheet number: 27
\deleted{Pour prouver} \inserted{Dire} que les produits finis indexés par l'ens. fini $I$ existent dans la catégorie $\mathcal A=\mathbb D(e)$ (disons) équivaut à dire que le foncteur diagonal
\[
\mathcal A\xrightarrow{\Delta}\mathcal A^I
\]
admet un adjoint à gauche. Or si
\[
f:I\longrightarrow e
\]
est la projection de $I$ (\deleted{\unsure{con}}considéré comme ens. ord. discret) sur \deleted{\unsure{le}} l'\unsure{en|ens} ordonné ponctuel $e$, le foncteur $\Delta$ s'identifie à $f^*:\mathbb D(e)\longrightarrow\mathbb D(I)$ (grâce à l'axiome $\unsure{\mathrm{D\bar er}|\mathrm{D\acute er}|\mathrm{Der}}_2$ sur le \unsure{comportement} des dérivateurs pour une somme \unsure{directe|discrète} de cat. d'indices).
Donc l'existence des produits revient à celle de l'adj. à gauche $f_*$ de $f^*$, \inserted{OK pour dérivateur g.} Et un hom. exact à gauche de dérivateurs à gauche, comme il commute aux $f_*$, commute aux produits finis. \unsure{Item|Itou} pour les sommes finies, et les dérivateurs gauches.

L'exemple précédent utilisait le fait que le dérivateur envisagé n'est pas pointé. Mais Même pour un dérivateur pointé \inserted{$\mathbb D$}, un morphisme $f^{\mathbb D}_*:\mathbb D_I\longrightarrow\mathbb D_{I'}$ n'est pas forcément exact à droite. Si on prend le dérivateur associé à une catégorie $\mathcal M$ quelc. (avec $\Sigma=\varnothing$),
\[
\mathbb D_{\mathcal M}(I)=\underline{\mathrm{Hom}}(I^\circ,\mathcal M)\overset{\text{déf}\unsure{\text{n}|\,}}{=}\mathcal M(I)
\]
alors le dérivateur induit
\[
\begin{aligned}
(\mathbb D_{\mathcal M})_I&=\big\{J\longmapsto\mathbb D_{\mathcal M}(I\times J)=\underline{\mathrm{Hom}}((I\times J)^\circ,\mathcal M))\\
&\hspace{6em}\simeq\underline{\mathrm{Hom}}(J^\circ,\underline{\mathrm{Hom}}(I^\circ,\mathcal M))
\end{aligned}
\]
\scanpage{29}
\mspage{27}
% archive sheet number: 28
s'identifie à
\[
\mathbb D_{\underline{\mathrm{Hom}}(I^{\unsure{\circ|'}},\mathcal M)}=\mathbb D_{\mathcal M(I)}
\]
i.e.
\[
(\mathbb D_{\mathcal M})_I\unsure{=|\simeq}\mathbb D_{\mathcal M(I)}
\]

\deleted{\unsure{Et si|Soit donc}
\[
I\xrightarrow{f}\overset{\unsure{\longleftrightarrow}}{I'}
\]
les foncteurs $f^*_{\mathbb D}$, $f_*^{\mathbb D}$, $f_!^{\mathbb D}$ entre $\mathbb D_I$ et $\mathbb D_{I'}$ s'interprètent comme des foncteurs entre $\mathbb D_{\mathcal M(I)}$ et $\mathbb D_{\mathcal M(I')}$.}

\uline{NB} Ici $\mathcal M$ et $\mathbb D_{\mathcal M}$ se déterminent mutuellement -- le foncteur
\[
\begin{aligned}
\mathcal M&\unsure{\mathord{-}|\longrightarrow}\mathbb D_{\mathcal M}\\
(\mathrm{Cat})&\longrightarrow\mathrm{DER}
\end{aligned}
\]
est 2-fidèle, i.e. si $\mathcal M,\mathcal M'$ sont \unsure{dans|deux} $\mathrm{Cat}$
\[
\underline{\mathrm{Hom}}(\mathcal M,\mathcal M')\xrightarrow{\approx}\mathrm{Hom}(\mathbb D_{\mathcal M},\mathbb D_{\mathcal M'})
\]
est une équivalence de catégories. D'autre part, on vérifie tout de suite que \deleted{$\unsure{f_I}$} \deleted{$\unsure{f}$} \inserted{$u$}$:\mathcal M\longrightarrow\mathcal M'$ est exact à g (resp. à droite) ssi. $\mathbb D_u:\mathbb D_{\mathcal M}\longrightarrow\mathbb D_{\mathcal M'}$ l'est.

Ainsi, pour voir si un \deleted{foncteur} \inserted{morphisme} de dérivateurs de la forme
\[
f^{\mathbb D}_*:\deleted{\unsure{\mathcal M_*}}\,
\underset{\substack{\unsure{\shortparallel|\simeq}\\\mathbb D_{\mathcal M(I)}}}{(\mathbb D_{\mathcal M}\deleted{\unsure{f}})_I}
\longrightarrow
\underset{\substack{\shortparallel\\\mathbb D_{\mathcal M(I')}}}{(\mathbb D_{\mathcal M})_{I'}}
\]
est exact à \uline{droite}, il revient au même de dire
\scanpage{30}
\mspage{28}
% archive sheet number: 29
que le \deleted{foncteur} pour les catégories fondamentales associées
\[
\mathcal M(I)\xrightarrow{f_*^{\mathcal M}}\mathcal M(I')
\]
est exact à droite. Or ici c'est un foncteur image directe de préfaisceaux ordinaires, et ce foncteur n'a aucune \unsure{envie|raison} d'être exact à \uline{droite} !
\margin{$!\ \shortparallel$}
\deleted{De fait, un \unsure{fonct}} \inserted{(Même si $\mathcal M$ est \uline{abélienne}, d'ailleurs !)}
Prenons p.ex. pour $f:I\hookrightarrow I'$ une imm. ouverte, avec $I'\simeq\Delta_1\times\Delta_1$, $I=I'\setminus\{\unsure{c|e}\}$ \inserted{objet final de $I'$}
\diagram{\begin{tikzcd}[ampersand replacement=\&,column sep=small,row sep=small]
\& c \& b\arrow[l] \& \& \& b\\[-1ex]
{I'\unsure{:|=}} \& a\arrow[u] \& \unsure{0|o}\arrow[u]\arrow[l] \& {I=} \& a \& 0\arrow[l]\arrow[u]
\end{tikzcd}}
Le foncteur $f_*$ est celui qui associe à un diagramme
\diagram{$F=\left\{\vcenter{\hbox{\begin{tikzcd}[ampersand replacement=\&,row sep=small]
\& F_b\arrow[d]\\
F_a\arrow[r] \& F_0
\end{tikzcd}}}\right\}$}
le carré cartésien complété
\diagram{$\bar F=\left\{\vcenter{\hbox{\begin{tikzcd}[ampersand replacement=\&,row sep=small]
F_a\times_{F_0}F_b\arrow[r]\arrow[d] \& F_b\arrow[d,"\beta"]\\
F_a\arrow[r,"\alpha"] \& F_0
\end{tikzcd}}}\right\}$}
\margin{Or si $F\mapsto\bar F$ était exact à dr., $F\mapsto F_a\times_{F_0}F_b$ \deleted{$\unsure{F_a\times}$} le serait aussi}

Or
\[
F_a\times_{F_0}F_b=K(F)=\mathrm{Ker}\big(F_a\times F_b\xrightarrow[\alpha\mathrm{pr}_1-\beta\mathrm{pr}_2]{}F_0\big)\,,
\]
en tant que foncteur additif en $F$, n'est pas exact à droite, p.ex. il ne transforme pas néc. épi en épi. Prenons par exemple
\diagram{$F'=\left\{\vcenter{\hbox{\begin{tikzcd}[ampersand replacement=\&,row sep=small]
\& F_b\arrow[d]\\
F_a\arrow[r] \& \deleted{\unsure{F_0}}0
\end{tikzcd}}}\right\}$}
et l'hom évident $F\longrightarrow F'$ qui est l'identité sur
\scanpage{31}
\mspage{29}
% archive sheet number: 30
$F_{\unsure{a|0}}$ et sur $F_b$. \inserted{C'est un épi, et} L'hom correspondant $\bar F\longrightarrow\overline{F'}$ est l'inclusion
\[
F_a\times_{F_{\unsure{0|*}}}F_b\longrightarrow F_a\times F_b
\]
qui n'est un épi que si $\alpha=\beta=0$.

\uline{Perplexité} : On devrait pouvoir définir une notion de dérivateur (additif) ``\uline{triangulé}'' (qui correspondrait à l'\unsure{intuition|intention} des cat. \deleted{\unsure{ab}} triangulées -- c'est un dérivateur satisfaisant à des axiomes assez forts), de sorte que pour les morphismes entre ceux-ci, l'exactitude à gauche et l'exactitude à droite soient \uline{\unsure{équivalentes|équivalents}}. \deleted{C'est peut-être}

\figure{\begin{tikzpicture}\draw[line width=2.8pt,line cap=round] (0,.38) -- (.23,0);\end{tikzpicture}}

\deleted{\unsure{Quest.}} \deleted{\unsure{Question de}}

``\uline{Théorèmes de changement de base}'' pour dérivateurs.

Soit $\mathbb D$ dérivateur à gauche (existence des $f_*$) et considérons
\[
f:I\longrightarrow J
\]
\inserted{dans $\mathrm{Diag}$}, soit $j_0\in\mathrm{Ob}\,J$, on se \deleted{\unsure{pr}}propose de ``calculer'' le foncteur
\[
F\longrightarrow f_*(F)_{j_0}=\alpha^*_{j_0}\big(f_*(F)\big)
\]
où
\[
\alpha_{j_0}:e\longrightarrow J
\]
est l'inclusion de la cat. ponctuelle $e$ dans $J$, de valeur \unsure{$j$.|$j_0$} \deleted{On voudrait trouver} \inserted{Considérons}
\scanpage{32}
\mspage{30}
% archive sheet number: 31
Considérons le diagramme
\diagram{\begin{tikzcd}[ampersand replacement=\&,row sep=large,column sep=large]
\& I_{/j_0}\arrow[r,"q"]\arrow[d,"f'"'] \& I\arrow[d,"f",""{name=side,left}]\\
{(1)} \& e\arrow[r,"\alpha_{j_0}","{\substack{\shortparallel\\p}}"',""{name=base,above}] \& J\arrow[from=side,to=base,dashed,bend right=12,"{\unsure{\gamma|\sigma}}"']
\end{tikzcd}}
où $I_{/j_0}$ est la catégorie \inserted{$I\times_J J_{/j_0}$} des couples
\[
\{(i,\varphi)\mid i\in\mathrm{Ob}\,I,\ \varphi:f(i)\longrightarrow j_0\}
\quad\deleted{\unsure{=I\times_J J_{/j_0}}}
\]
et $q$ est le foncteur $(i,\varphi)\longmapsto i$. Ce diagramme n'est \uline{pas} commutatif, mais il y a un foncteur canonique
\[
fq\xrightarrow{\gamma}pf'
\]
où $fq$ est le foncteur $(i,\varphi)\mapsto f(i)$ et $pf'$ le foncteur constant $(i,\varphi)\longmapsto j_0$ et où $\gamma(i,\varphi):f(i)\longrightarrow j_0$ est défini par $\underline{\gamma(i,\varphi)=\varphi}$.

Plus généralement, considérons un diagramme dans $\mathrm{Diag}$
\diagram{\begin{tikzcd}[ampersand replacement=\&,column sep=large]
\& I'_{\deleted{\unsure{/j_0}}}\arrow[r,"q"]\arrow[d,"f'"'] \& I\arrow[d,"f",""{name=side,left}]\\
{(2)} \& J'\arrow[r,"p",""{name=base,above}] \& J\arrow[from=side,to=base,bend right=8,"\gamma"'] \& {\text{avec }\gamma:fq\to pf'.}
\end{tikzcd}}
On définit alors un \uline{morphisme de chgt de base pour les images directes} relatif à ce diagramme
\[
\boxed{p^*f_*\xrightarrow{c_\gamma}f'_*q^*}
\]
\scanpage{33}
\mspage{31}
% archive sheet number: 32
de deux façons, qui donnent le $=$ résultat \deleted{\unsure{évident|évid.}} (histoire générale de foncteurs adjoints, à expliciter une bonne fois, en \unsure{termes|terme} d'un diagramme de foncteurs
\diagram{\begin{tikzcd}[ampersand replacement=\&]
 X' \arrow[d,leftrightarrow,"{f'^{*}_{*}}"',""{name=vert,right}] \& X \arrow[l,"{q^*}"',""{name=top,below}]\\
 Y' \arrow[from=vert,to=top,bend right=17,shorten >=1pt] \& Y \arrow[l,"{p^*}"'] \arrow[u,leftrightarrow,"{\unsure{F|f}^*}"']
\end{tikzcd}}
en supposant que les foncteurs adj. à g. $f'_*,f_*$ des deux foncteurs verticaux existent).

Un cas intéressant est celui\deleted{\unsure{ci}} où
\[
 pf'=fq,\qquad \unsure{\gamma|\sigma}=\mathrm{id}_{fq}
\]
et où on \deleted{\unsure{peut}} pose $c_\gamma=c$.

\diagram{\begin{tikzcd}[ampersand replacement=\&,column sep=small]
 \text{Ici} \& p^*f_* \arrow[r,"{?}"] \& f'_*q^*\\
 \text{revient à} \& f_* \arrow[r,"{?}"] \& \underbrace{p_*f'_*}_{\deleted{\unsure{f_*|f}}}q^*\simeq\tikz[baseline=(under.base)]{\node[inner sep=0pt] (under) {$(pf')_*$}; \draw ([yshift=1pt]under.south west) -- ([yshift=-1pt]under.south west) -- ([yshift=-1pt]under.south east) -- ([yshift=1pt]under.south east);}q^*\\
 \& \& (fq)_*q^*\simeq(f_*q_*)q^*=f_*(q_*q^*) \arrow[u,"{\gamma_* *q^*}"',"{\unsure{\prime\prime|\shortparallel}}"]\\
 \& \deleted{f_*} \arrow[ur,bend right=10,"{\deleted{f_* *\alpha}}",postaction={decorate,decoration={zigzag,segment length=7pt,amplitude=2pt,post length=15pt}}] \&
\end{tikzcd}}
Or on a $\mathrm{id}\xrightarrow{\alpha}q_*q^*$ (adjonction)

d'où $f_*\xrightarrow{f_* *\alpha}f_*(q_*q^*)$

et en composant avec $\gamma*q^*$, on trouve ce qu'on veut. Et on peut au lieu de ceci chercher
\diagram{\begin{tikzcd}[ampersand replacement=\&,column sep=small]
 f'^*(p^*f_*) \arrow[r,"{?}"] \& q^*\\[-1ex]
 \overset{\unsure{\shortparallel|\prime\prime}}{(pf')^*}f_* \arrow[d,"{\gamma^* * f_*}"] \& \\
 (fq)^*f_*\simeq q^*(f^*\deleted{\unsure{_*}}f_*) \arrow[r,"{q^* *\gamma'}"] \& q^* \& {\text{où }\tikz[baseline=(letter.base)]{\node[inner sep=0pt] (letter) {$\gamma'$}; \node[inner sep=0pt,anchor=north west] (adj) at ([xshift=1pt,yshift=-3pt]letter.south east) {$\scriptstyle\text{adj.}$}; \draw[->] (letter.south west) to[bend right=25] (adj.west);}\colon f^*f_*\longrightarrow\mathrm{id}}
\end{tikzcd}}
